% THE ARTICLE'S SUBSET of si_tables/external_baselines_rows.tex, split off 2026-08-08.
% That file stays the full transcription and is what the Supporting Information prints
% (Table \ref{tab:si-baselines-full}); this one is the same rows with two blocks' worth
% removed so the article's copy can be set at \footnotesize instead of \scriptsize.
%
% WHAT IS NOT HERE, AND WHY EACH WENT.  Three referees and the supervisor all named this
% table's type size; at 8 pt it filled 82% of a float page carrying no body text, and the
% only way to 10 pt was to take rows out.
%   - the third block, six rows quoted AS PUBLISHED on other corpora, in mol/L, with a
%     dash in three of the four data columns.  No value in it is computed here and no
%     conversion of ours stands behind any of them, so a reader cannot check them against
%     anything else in the table anyway; the claim the article makes against them prints
%     their range in the sentence that makes it (\S\ref{sec:discussion}).
%   - the two RANDOM-PAIR rows.  They are a split control -- the same two arms on a split
%     where the solute is seen in training -- and the article makes no claim from them.
%   - the note block, 489 words in twenty lines, which is the whole of it in the SI copy.
%   - THE `Test set' COLUMN, dropped 2026-08-08.  It printed `solute scaffold, this corpus'
%     five times and `by-solute, this corpus' four times, which is what the two italic block
%     heads already say -- so it carried no information at all, and it was the widest thing
%     between the model names and the right margin.  Its 116 pt is what wrapped every model
%     cell onto two or three lines, and the block contrast is what this table exists to show.
%     The one thing the column said that a head did not was the word `by-solute', which is
%     now in the third head.  Do NOT restore the column; if a future block needs a set that
%     no head states, put it in that block's head.
% NO VALUE AND NO ROW ORDER IS CHANGED between the two files.  The one textual
% difference in a row is on the surviving FastSolv row: the full file keys it \tnk{a} to the note
% that says what "re-scored" means, and with no note block here the word is set in the
% row instead.  If the numbers move, re-run the generator, re-apply the four typesetting
% edits listed in the full file, and re-cut this one from it -- then drop the Test-set column
% and the two blocks above, which is the whole of the difference.
% THE COLUMN HEADS ARE STACKED HERE AND NOT IN THE FULL FILE, and that is arithmetic rather
% than taste.  The L{} box is a fraction of \linewidth, so it does not shrink when the
% type does, while the four numeric columns are set to their widest cell and do.  At 8 pt the
% whole thing measured 485 pt inside a 504.5 pt full-width float; at 10 pt the same specifier
% measures 533 pt and printed 30 pt into the right margin (one Overfull \hbox).  Two heads were
% wider than any number under them -- `n (ln x2/log10 S)' by 10 pt and `RMSE log10 S' by 14 pt --
% so breaking all four over two lines returns 24 pt at no cost to any cell.
% THE MEASURE, RE-DONE FOR FIVE COLUMNS (2026-08-08).  \textwidth in this document's table* is
% 504.47 pt, measured with \typeout and not assumed.  The four numeric columns cost 217.2 pt
% together (set by `(ln x2/log10 S)' at 62.7, `-0.03 +- 0.04' at 55.6 and two `1.70 +- 0.03'
% at 47.8), and five columns with @{} at both ends cost 2*\tabcolsep*4 in gaps.  The longest
% model name -- `TGNN, COSMO-SAC, reference sigma* substituted at evaluation' -- has a natural
% width of 268.4 pt, so at \tabcolsep 3 pt the model column could be at most 263.3 pt and that
% name would still wrap.  \tabcolsep 2 pt (in the article, beside the \input) buys the 8 pt that
% closes it: 269.9 + 217.2 + 16 = 503.1 pt, which is 1.4 pt inside the measure, and every one of
% the nine names then sets on one line.  Both margins are thin and NOTHING here is stretchy, so
% anything that widens a cell or a block head must be re-measured against 504.47 pt --
% \sbox the tabular and \typeout its \wd; the build does not reliably complain, because an
% over-wide tabular in a centred float sticks out on both sides without an Overfull \hbox.
% The widest block head is the third at 486.7 pt, which is inside the tabular and not binding.
\begin{tabular}{@{}L{0.535\linewidth}rrrr@{}}
\toprule
Model & \shortstack[r]{$n$\\($\lnx$/$\log_{10}S$)} & \shortstack[r]{MAE\\$\lnx$} &
\shortstack[r]{RMSE\\$\log_{10}S$} & \shortstack[r]{$R^2$\\$\lnx$} \\
\midrule
\multicolumn{5}{@{}l@{}}{\emph{This work, solute-scaffold split; mean $\pm$ s.d. over seeds, three unless the row says otherwise}}\\
DirectGNN, same encoder, no closure \emph{(control)} & $5608$ / $5440$ & $1.70\pm0.03$ & $1.00\pm0.02$ & $0.42\pm0.03$ \\
TGNN, NRTL closure (physics) & $5608$ / $5440$ & $1.79\pm0.07$ & $1.54\pm0.07$ & $0.34\pm0.03$ \\
TGNN, COSMO-SAC closure, supervised $\hat\sigma$ (physics) & $5608$ / $5440$ & $1.85\pm0.05$ & $1.39\pm0.19$ & $0.33\pm0.05$ \\
TGNN, COSMO-SAC, reference $\sigma^\star$ substituted at evaluation & $5608$ / $5440$ & $2.25\pm0.02$ & $1.60\pm0.14$ & $-0.03\pm0.04$ \\
\addlinespace[2pt]
\multicolumn{5}{@{}l@{}}{\emph{Scale reference on the same solute-scaffold rows}}\\
predict the test-set mean of $\lnx$ & $5608$ / $5440$ & $2.32$ & $1.33$ & $0.00$ \\
\addlinespace[2pt]
\multicolumn{5}{@{}l@{}}{\emph{External models re-run on this corpus on a by-solute split, one run each---not a ranking against the block above}}\\
FastSolv, retrained on this corpus, re-scored & $10287$ & $1.44$ & $0.84$ & $0.55$ \\
SolProp cycle, released model $+$ 3-parameter recalibration & $10287$ & $2.34$ & $1.34$ & $-0.18$ \\
SolProp cycle, released model, zero-shot & $10287$ & $2.36$ & $1.87$ & $-0.47$ \\
SolProp encoder retrained on $\lnx$, cycle discarded & $10570$ / $10287$ & $1.62$ & $1.07$ & $0.39$ \\
\bottomrule
\end{tabular}
